% Work-MAT | 2. rocnik | hodina 16
\documentclass[11pt,a4paper]{article}
\input{../../sablony/pripravy/priprava.sty}
\usepackage{tikz}
\usetikzlibrary{arrows.meta}
\nastavPripravu{2. ročník}{Exponenciální rovnice}{16}{Exponenciální rovnice IV: substituce}
\begin{document}
\poznamka{Vytýkání z minulé hodiny je často nejrychlejší, dnes však žáky vedeme k pojmenování opakujícího se výrazu $t=a^x$. Orientačně 8 min motivace, 14 min modelový postup, 12 min samostatná práce, 11 min rozbor a shrnutí.}
\blok{Z opakující se mocniny uděláme novou neznámou}
V rovnici $2\cdot3^x+3^x=27$ se opakuje výraz $3^x$. Je možné jej vytknout, ale také dočasně označit písmenem $t$. Tím se původní exponenciální rovnice změní na známou lineární rovnici.
\begin{definiceBox}[title={Definice - Substituce exponenciálního výrazu}]
\textbf{Substituce} $t=a^x$ znamená, že opakovaně vystupující výraz $a^x$ dočasně nahradíme novou neznámou $t$. Pro $a>0$ platí $t>0$. Po výpočtu $t$ se musíme vrátit k rovnici pro $x$.
\end{definiceBox}
\begin{prikladReseniBox}[title={Příklad 1}]
Vyřeš $2\cdot3^x+3^x=27$ pomocí substituce.
\tcblower
\[\begin{aligned}t&=3^x,\qquad t>0,\\2t+t&=27\\3t&=27\\t&=9.\end{aligned}\]
Zpětná substituce:
\[\begin{aligned}3^x&=9\\3^x&=3^2\\x&=2.\end{aligned}\]
Zkouška: $2\cdot9+9=27$. $K=\{2\}$.
\end{prikladReseniBox}
\blok{Kdy je pomocná hodnota přípustná?}
Rovnice pro $t$ může dát i záporný výsledek. Ten ale nelze přijmout, protože exponenciální výraz $a^x$ je vždy kladný.
\begin{prikladReseniBox}[title={Příklad 2}]
Vyřeš $4\cdot2^x-20=12$.
\tcblower
\[\begin{aligned}t&=2^x,\quad t>0,\\4t-20&=12\\4t&=32\\t&=8.\end{aligned}\]
Návrat: $2^x=8=2^3$, tedy $x=3$. Zkouška: $4\cdot8-20=12$.
\end{prikladReseniBox}
\begin{prikladReseniBox}[title={Příklad 3}]
Vyřeš $2\cdot5^x+7=5$.
\tcblower
Položíme $t=5^x>0$. Dostaneme $2t+7=5$, tedy $t=-1$. To odporuje $t>0$, proto původní rovnice nemá reálné řešení: $K=\varnothing$.
\end{prikladReseniBox}
\poznamka{Nejdůležitější čtyři kroky napiš na tabuli: 1. substituce a podmínka; 2. nová rovnice; 3. řešení a kontrola $t>0$; 4. návrat k $x$ a zkouška. Žáci často skončí výsledkem $t=8$ a zapomenou na původní otázku.}
\begin{zadaniReseniBox}[title={Samostatná práce 1}]
Použij substituci $t=a^x>0$.
\begin{enumerate}[label=\alph*)]
\item $3\cdot2^x+2^x=32$;
\item $2\cdot5^x-10=40$;
\item $5\cdot3^x+5=0$.
\end{enumerate}
\tcblower
\begin{enumerate}[label=\alph*)]
\item $4t=32$, $t=8$, $2^x=8$, tedy $x=3$.
\item $2t=50$, $t=25$, $5^x=25$, tedy $x=2$.
\item $t=-1$ odporuje $3^x>0$, proto $K=\varnothing$.
\end{enumerate}
\end{zadaniReseniBox}
\begin{zadaniReseniBox}[title={Rozšíření 1}]
Vyřeš $3\cdot2^{x+1}-2^{x-1}=44$ pomocí $t=2^x$.
\tcblower
$3\cdot2t-\frac12t=44$, tedy $\frac{11}{2}t=44$, $t=8>0$, takže $2^x=8$ a $x=3$.
\end{zadaniReseniBox}
\begin{zadaniReseniBox}[title={Domácí úkol}]
\begin{enumerate}[label=\alph*)]
\item Substitucí vyřeš $5\cdot2^x-8=32$.
\item Substitucí vyřeš $3\cdot4^x+5=8$.
\item Stručně vysvětli, proč substituce $t=2^x$ vyžaduje $t>0$.
\end{enumerate}
\tcblower
\begin{enumerate}[label=\alph*)]
\item $t=8$, $2^x=8$, $x=3$.
\item $3t=3$, $t=1$, $4^x=1$, $x=0$.
\item Každá reálná mocnina kladného základu $2$ je kladná.
\end{enumerate}
\end{zadaniReseniBox}
\end{document}
